\documentclass[11pt]{article}
\usepackage{amsmath,amssymb,amsthm}
\newtheorem{lemma}{Lemma}
\newtheorem{theorem}{Theorem}
\begin{document}
\title{Higher-level interference in shared-residue Haar cardinalization}
\author{}
\date{}
\maketitle

\paragraph{Status and scope.}
This is a complete written obstruction to one proposed decoder, with exact
finite routing and four-wire checks. It is not Lean checked and is not a
lower bound for the original HeadComplexity model. The restrictions are
one fixed residue pool, one exterior Cauchy kernel per query, pure input
Haar profiles, and correctness on every setting of the butterfly network.
Target-dependent residues, several kernels in one query, different input
profiles, and a selected subfamily whose later switches serve as auxiliary
variables are not covered.

Use the butterfly notation of \texttt{butterfly-wavelet-span.tex}.
There are $H=2^d$ wires. Fix real residues $r_y$, poles
$0\leq\xi_y\leq a$, and a positive integer $p$. Write
\[
 w_y(t)=\frac{r_y}{(t-\xi_y)^p},\qquad t\geq b>a.
\]
For a fixed level $k$ and input block $g$, queries have scores
\[
 F_s(\pi)=\beta_s+\lambda_s\sum_h u_{k,g}(h)w_{\pi(h)}(t_s),
 \qquad 0\leq s<2^k,
\]
where $t_s\geq b$, $\lambda_s\ne0$, and $\beta_s$ are arbitrary.
The desired strict sign is $1-2\theta_{k,s,g}$, where $\theta$ is the
corresponding switch bit.

\begin{lemma}[Independent later routing]
Choose one representative $y_a$ from every output residue class
$a\pmod {2^{k+1}}$. For every choice of the $2^k$ layer-$k$ switch bits
in block $g$, there is a butterfly setting such that the input
\[
 h=(g\,2^{k+1})+b'2^k+s
\]
is sent to $y_{(b'\mathbin{\oplus}\theta_{k,s,g})2^k+s}$ for every
$s<2^k$ and $b'\in\{0,1\}$.
\end{lemma}
\begin{proof}
Set every earlier switch to the identity. Set the chosen layer-$k$
switches as required. The tracked wires then have distinct lower
$k+1$ tuples. Later layers preserve those tuples, so tracked wires
never share a later switch. At each later layer, set the switch on
each tracked path to give its next bit the value of its chosen output
representative. These prescriptions never conflict, and after the last
layer every tracked wire has the prescribed output.
\end{proof}

\begin{theorem}[A uniform bound on decodable prefixes]
If the displayed queries decode their respective layer-$k$ switch bits
on every butterfly setting, then
\[
 2^k<(p+1)\left[1+\left(\frac b{b-a}\right)^{4p}\right].
 \tag{1}
\]
In particular, for the simple Cauchy kernel $p=1$ and
$a/b<1-3^{-1/4}$, no level $k\geq3$ admits this decoder.
\end{theorem}
\begin{proof}
Fix arbitrary representatives as in the lemma and put
\[
 \Delta_s(t)=w_{y_s}(t)-w_{y_{2^k+s}}(t).
\]
For the settings constructed by the lemma, the score before its bias
and scale is exactly
\[
 \sum_{j<2^k}(1-2\theta_{k,j,g})\Delta_j(t).
\]
All signs in this sum vary independently. Correctness of query $s$
for every such choice forces
\[
 |\Delta_s(t_s)|>\sum_{j\ne s}|\Delta_j(t_s)|.
 \tag{2}
\]
Indeed the target term must overcome both choices of its own sign and
the worst choices of all other signs; an arbitrary bias cannot weaken
this necessary inequality. In particular each diagonal term is nonzero.

Let $R=2^k$. Write $\Delta_j=N_j/D_j$, where
\[
 D_j(t)=(t-\xi_{y_j})^p(t-\xi_{y_{R+j}})^p
\]
is positive and $N_j$ is a polynomial of degree at most $p$.
The matrix
\[
 M_{ij}=\frac{N_j(t_i)}{N_i(t_i)}
\]
has diagonal entries one and rank at most $p+1$. Since
$\Delta_j(t_i)/\Delta_i(t_i)=M_{ij}D_i(t_i)/D_j(t_i)$,
(2) and the pole interval give
\[
 \sum_{j\ne i}|M_{ij}|<L,\qquad
 L=\left(\frac b{b-a}\right)^{2p}.
\]
Here $D_j(t_i)/D_i(t_i)\leq(t_i/(t_i-a))^{2p}\leq L$.
Consequently $\|M\|_F^2<R(1+L^2)$. For any real matrix,
$|\operatorname{tr}M|\leq\sqrt{\operatorname{rank}M}\|M\|_F$:
apply Cauchy--Schwarz to its singular values and use that the absolute
trace is at most their sum. Thus
\[
 R^2<(p+1)R(1+L^2),
\]
which is (1). In the final case its right side is less than eight.
\end{proof}

\paragraph{Why the interpolation count does not settle the route.}
The same-level Haar coefficients
\[
 W_{k,s}(t)=\sum_{y\bmod 2^k=s}(-1)^{y_k}w_y(t)
\]
can be constrained to vanish at query times belonging to other prefixes.
For $k\leq\frac12\log_2H$, the total number of such linear conditions
is only $O(H)$. However these conditions control averages over the later
output bits. The independent-routing lemma lets those bits select
individual representatives instead. Inequality (1) shows that no choice
of the $H$ residues can make their unaveraged contributions harmless on
all network settings at a growing number of levels. This identifies the
higher-level interference gap quantitatively; it does not rule out using
later switches as target-dependent auxiliary choices.

\paragraph{A nontrivial four-wire exact example.}
The bound is not a claim that every multilevel decoder fails. Let
$\epsilon=1/10000$, $T=99/100$, $\xi_j=j\epsilon$ for $0\leq j<4$, and
\[
 r_0=100/\epsilon,\quad r_1=0,\quad
 r_2=r_0(T-2\epsilon)/T,\quad r_3=-1.
\]
Use time one for both level-zero queries and for the level-one prefix-zero
query, and use time $T$ for the level-one prefix-one query. All four
switch signs are decoded on all sixteen network settings.
For level zero, both even-indexed wire weights are positive and both
odd-indexed weights are nonpositive, so every routed difference has the
required strict sign. At level one the two gaps are
\[
 \Delta_0(t)=\frac{2\epsilon r_0(t-T)}{Tt(t-2\epsilon)},
 \qquad \Delta_1(t)=\frac1{t-3\epsilon}.
\]
At $t=T$ the first is zero and the second positive. At $t=1$ the first
is larger than the second, as follows directly from the displayed
rational values. This proves both independent level-one signs.

The companion checker verifies all 64 strict scores of this example
using exact rational arithmetic. It also checks 13,862 instances of the
routing lemma, exhausting all representative choices, switch bits, and
input blocks for $d\leq4$. Exploratory floating-point infeasibility on
larger examples is not used in the theorem.
\end{document}
